\chapter{Подготовка и предобработка данных}

\section{Разведочный анализ и структура данных}

Исходная выборка насчитывает 2500 записей и 10 атрибутов. Первичный анализ выявил неоднородность типов: 3 числовых признака, 6 категориальных и 1 числовой целевой показатель.

\begin{table}[H]
\centering
\small
\caption{Статистика пропущенных значений в признаках}
\begin{tabular}{|l|c|c|l|}
\hline
\textbf{Признак} & \textbf{Пропусков (шт.)} & \textbf{Доля (\%)} & \textbf{Метод заполнения} \\
\hline
\texttt{bedrooms} & 102 & 4.1\% & Медиана (3.0) \\
\texttt{bathrooms} & 119 & 4.8\% & Медиана (2.0) \\
\texttt{level} & 166 & 6.6\% & Мода (Ground) \\
\texttt{finishing} & 101 & 4.0\% & Мода (Finished) \\
\hline
\end{tabular}
\end{table}

Пропуски в объявлениях вызваны человеческим фактором при заполнении форм риелторами. Пропущенные непрерывные значения заполнены медианой, так как она устойчива к отдельным всплескам и асимметрии распределения. Категориальные признаки дополнены наиболее частым классом (модой).

\begin{lstlisting}[language=Python]
# Заполнение пропусков устойчивыми статистиками
df_clean = df.copy()
for col in ['bedrooms', 'bathrooms']:
    df_clean[col] = df_clean[col].fillna(df_clean[col].median())

for col in ['level', 'finishing']:
    df_clean[col] = df_clean[col].fillna(df_clean[col].mode()[0])
\end{lstlisting}

\section{Кодирование категориальных признаков}

В строгом соответствии с пунктом 1.2 методических указаний категориальные признаки преобразованы с помощью \texttt{LabelEncoder}. Каждому уникальному текстовому значению ставится в соответствие целочисленный индекс от $0$ до $K-1$:

\begin{lstlisting}[language=Python]
from sklearn.preprocessing import LabelEncoder

categorical_features = ['type', 'level', 'city', 'finishing', 'furnished', 'payment_method']
label_encoders = {}
for col in categorical_features:
    le = LabelEncoder()
    df_clean[col] = le.fit_transform(df_clean[col].astype(str))
    label_encoders[col] = le
\end{lstlisting}

\section{Разбиение выборки и масштабирование признаков}

Выборка разделена на обучающую (80\%, 2000 объектов) и контрольную (20\%, 500 объектов) при фиксированном значении генератора случайных чисел \texttt{random\_state=42}.

\begin{lstlisting}[language=Python]
from sklearn.model_selection import train_test_split
from sklearn.preprocessing import StandardScaler

X = df_clean.drop(columns=['price'])
y = df_clean['price']

X_train, X_test, y_train, y_test = train_test_split(
    X, y, test_size=0.20, random_state=42
)

scaler = StandardScaler()
X_train_scaled = scaler.fit_transform(X_train)
X_test_scaled = scaler.transform(X_test)
\end{lstlisting}

Масштабирование признаков через \texttt{StandardScaler} вычисляет центрирование и нормировку:
\begin{equation}
z_{ij} = \frac{x_{ij} - \mu_j}{\sigma_j}
\end{equation}
где $\mu_j$ и $\sigma_j$ — выборочные среднее и стандартное отклонение $j$-го признака на обучающей подвыборке. Параметры скейлера оцениваются исключительно по обучающей выборке, чтобы избежать утечки информации из теста.

\section{Выводы по разделу 1}

Проведенный этап устранил неполноту исходных табличных данных и нормализовал признаковое пространство. Устранение дисбаланса масштабов (когда площадь измерялась сотнями, а число комнат — единицами) является обязательным условием перед обучением регрессий с $L_1$ и $L_2$ штрафами. Без стандартизации штраф действовал бы преимущественно на признаки с искусственно завышенными коэффициентами.
